%% SECTION 7: convexity-preserving column-based extension (FW / origin methods).
%% Added 2026-07-18. All reported numbers trace to
%% compressed_TAP_repro_package_v1_1_FW_origin/fw_origin_extension/results/
%% (reference_sioux/FW_controlled.csv, reference_sioux/ORIGIN_controlled.csv,
%%  reference_sioux/cpp_kernel_sioux.csv, reference_chicago/cpp_kernel_chicago.csv).

\section{A Convexity-Preserving Column-Based Extension: Frank--Wolfe and Origin-Decomposed Methods}
\label{sec:fw_origin}

The grouped representation of Section~\ref{sec:grouped} was introduced as a
constraint-compatible refinement of the compressed formulation solved by the AL
method. In this section we show that the same construction carries a stronger
and broader interpretation: it is a \emph{convexity-preserving reformulation}
of the path-based problem, not merely a heuristic grouping device, and it is
therefore directly compatible with classical feasible-direction methods ---
Frank--Wolfe (FW) and origin-decomposed conditional gradient schemes --- whose
theory rests on convex combinations of feasible assignments. The central
observation is that each latent route family is a convex combination of
feasible paths, the link-flow mapping remains linear, and the Beckmann
objective remains convex; consequently every statement of the classical convex
theory survives compression unchanged.

\subsection{Latent Route Families as Convex Combinations}
\label{sec:latent-convex}

Consider OD pair $i$ and let $b_p \in \Re^m$ denote the $p$th row of the
path-link incidence matrix $B$, written as a column vector. As in
Section~\ref{sec:grouped}, partition the minor paths of OD pair $i$ into
groups $G_{ik}$, $k = 1, \dots, r_i$, with fixed shares satisfying
\eqref{eq:group-shares}:
$\pi_{ikp} \ge 0$ and $\sum_{p \in G_{ik}} \pi_{ikp} = 1$.
The aggregate column of family $(i,k)$, cf.\ \eqref{eq:aggregate-column}, is
\begin{equation}
\label{eq:family-hull}
\bar b_{ik} = \sum_{p \in G_{ik}} \pi_{ikp}\, b_p
\;\in\;
\operatorname{conv}\bigl\{\, b_p : p \in G_{ik} \,\bigr\}.
\end{equation}
A latent route family is therefore not necessarily one physical path; it is a
feasible convex combination of physical paths of the same OD pair. This is the
mathematically precise meaning of the latent route family.

Let $y_i \ge 0$ collect the flows on the explicit major paths of OD pair $i$
and $s_{ik} \ge 0$ the flow assigned to family $G_{ik}$, subject to
\begin{equation}
\label{eq:z-simplex}
\sum_{p \in \mathcal M_i} y_{ip} + \sum_{k=1}^{r_i} s_{ik} = d_i ,
\end{equation}
where $\mathcal M_i$ is the major-path set of OD pair $i$. The implied
minor-path flows are $x_p = \pi_{ikp}\, s_{ik}$ for $p \in G_{ik}$, cf.\
\eqref{eq:group-flow}. Stacking $z_i = (y_i, s_i)$ and $z = (z_1, \dots,
z_\ell)$, the reconstruction of the full path-flow vector is linear,
\begin{equation}
\label{eq:G-reconstruction}
x = G z ,
\end{equation}
where $G$ contains an identity column for each explicit major path and a
nonnegative convex-combination column (the shares $\pi_{ikp}$) for each latent
route family. The following proposition records the properties that make this
representation exact convex programming rather than approximation heuristics.

\begin{proposition}[Convexity preservation of latent path-space compression]
\label{prop:convexity-preservation}
Let each compressed column of OD pair $i$ be either a feasible path-incidence
vector (a major path) or a convex combination \eqref{eq:family-hull} of
feasible path-incidence vectors of the same OD pair. Then:
\begin{itemize}
\item[(a)] every compressed solution $z \ge 0$ satisfying
\eqref{eq:z-simplex} reconstructs, via \eqref{eq:G-reconstruction}, a
nonnegative path-flow vector satisfying exact OD conservation, so the
represented path-flow set $\mathcal X_G = \{\, Gz : z \in \mathcal Z \,\}$,
with $\mathcal Z = \prod_i \{ z_i \ge 0 : \mathbf 1' z_i = d_i \}$, is a
convex subset of the original feasible path-flow set $\mathcal X$ of problem
\eqref{eq:original_ta};
\item[(b)] if each link travel-time function is nondecreasing, the compressed
problem
\begin{equation}
\label{eq:compressed-fw}
\min_{z \in \mathcal Z} \; f\bigl(\bar B' z + v_0\bigr),
\qquad \bar B = G' B ,
\end{equation}
is a convex program over a product of OD simplices.
\end{itemize}
\end{proposition}

\begin{proof}
(a) Nonnegativity: $z \ge 0$ and $G \ge 0$ entrywise imply $x = Gz \ge 0$; in
particular $x_p = \pi_{ikp} s_{ik} \ge 0$ for every minor path. Conservation:
for OD pair $i$,
$\sum_{p \in G_{ik}} x_p = s_{ik} \sum_{p \in G_{ik}} \pi_{ikp} = s_{ik}$,
so the total path flow of OD pair $i$ equals
$\sum_{p} y_{ip} + \sum_k s_{ik} = d_i$ by \eqref{eq:z-simplex}. Hence $Gz \in
\mathcal X$. Each $\mathcal Z_i$ is a simplex, so $\mathcal Z$ is convex, and
$\mathcal X_G$ is the linear image of a convex set, hence convex.
(b) The link-flow map $v = B'x + v_0 = B'Gz + v_0 = \bar B' z + v_0$ is affine
in $z$, and $f$ is convex when the link travel-time functions are
nondecreasing; therefore $z \mapsto f(\bar B' z + v_0)$ is convex, and the
constraint set is a product of simplices.
\end{proof}

The method thus solves the original convex assignment problem over a
compressed convex \emph{inner approximation} $\mathcal X_G \subseteq \mathcal
X$ of the path-flow feasible region.

\subsection{Two Latent Solvers: L\,$+$\,FW and L\,$+$\,PG}
\label{sec:fw-interpretation}

Because the feasible set of \eqref{eq:compressed-fw} is a product of OD
simplices, any solver for smooth convex optimization over this constraint
geometry applies to the latent representation unchanged. We state the two
that we implement and certify.\footnote{Reference implementations:
\texttt{fw\_origin\_compression.py} (latent FW with exact BPR line search and
the origin-block variant) and \texttt{latent\_projected\_gradient.py} (latent
projected gradient, about $90$ lines on top of the shared representation
layer), both in \texttt{compressed\_TAP\_repro\_package\_v3\_1\_LPG/repro/fw/src/};
a C++ Barzilai--Borwein projected-gradient kernel operating on the exported
\texttt{group.csv} representation is in \texttt{repro/cpp/tapkernel.cpp}.
Both Python solvers share the atom set, the latent gradient, the simplex
feasible set, and the FW duality-gap stopping certificate.}

\emph{Latent Frank--Wolfe} (L\,$+$\,FW). At iterate $z^k$ with link flows
$v^k = \bar B' z^k + v_0$, the linearization step
\begin{equation*}
w^k \in \arg\min_{w \in \mathcal Z}\; \nabla_z f\bigl(\bar B' z^k + v_0\bigr)' w
\end{equation*}
decomposes by OD pair into the standard minimum-cost column selection --- the
same linear minimization oracle as in classical path-based FW, now over major
columns and family columns priced by $\bar b_{ik}' t(v^k)$. The update
$z^{k+1} = (1 - \alpha_k) z^k + \alpha_k w^k$, with $\alpha_k$ from an exact
BPR line search on $[0,1]$, stays in $\mathcal Z$ by convexity, and in
original path-flow coordinates
\begin{equation*}
x^{k+1} = G z^{k+1} = (1 - \alpha_k)\, x^k + \alpha_k\, G w^k ,
\end{equation*}
so every iterate remains a convex combination of feasible path-flow
assignments of the original problem.

\emph{Latent projected gradient} (L\,$+$\,PG). The same latent gradient
$g^k = \bar B\, t(\bar B' z^k + v_0)$ instead drives a projected-gradient
update on the product of OD simplices,
\begin{equation*}
z^{k+1} = \Pi_{\mathcal Z}\bigl( z^k - \gamma_k\, g^k \bigr),
\end{equation*}
where $\Pi_{\mathcal Z}$ is the Euclidean projection (computed independently
per OD simplex in $O(K \log K)$), and $\gamma_k$ is a safeguarded
Barzilai--Borwein trial step with Armijo backtracking on the projected
direction. We deliberately do not call this method a reduced-gradient
algorithm: there is no basis, no elimination, and no reduced costs --- only
projection onto the same demand simplices that L\,$+$\,FW uses. Both solvers
stop on the identical FW duality gap, computed from the same linear
minimization oracle, so accuracy is matched by construction.
Feasibility of every L\,$+$\,PG iterate follows from
Proposition~\ref{prop:convexity-preservation}(a) exactly as for L\,$+$\,FW:
$z^{k+1} \in \mathcal Z$ implies $x^{k+1} = G z^{k+1} \in \mathcal X$. The representation does not introduce a
signed decoder or a nonconvex latent model: it replaces selected path columns
by feasible convex combinations of path columns, preserving the convex
feasible region, the convex Beckmann objective, and the classical
convex-combination structure of FW iterations; cf.\
\cite{bertsekas2009convex}. The same argument applies verbatim to
origin-decomposed conditional gradient schemes, since the family columns are
OD-local and hence origin-local: compression acts within each origin block and
exact conservation is retained per block.

\subsection{Exact Pricing and Promotion Preserve Convexity}
\label{sec:promotion-convex}

Suppose periodic exact pricing (a shortest-path computation at the current
link times) identifies a path $p^\star$ absent from the representation, and
$p^\star$ is promoted to the major set: its incidence vector is appended as a
new identity column of $G$. Each promotion only enlarges the represented
region, giving the nested sequence
\begin{equation*}
\mathcal X_{G^{(0)}} \subseteq \mathcal X_{G^{(1)}} \subseteq
\mathcal X_{G^{(2)}} \subseteq \cdots \subseteq \mathcal X ,
\end{equation*}
with every intermediate problem convex by
Proposition~\ref{prop:convexity-preservation}. Consequently the restricted
optimal objectives are monotonically nonincreasing,
$f^{(0)} \ge f^{(1)} \ge \cdots \ge f^\star$, and if sufficiently many actual
path columns are eventually promoted, the compressed feasible region
approaches the full path-flow feasible region. This is the standard
column-generation guarantee, inherited without modification because promotion
never leaves the convex-combination framework.

An important distinction from the signed SVD basis of
Section~\ref{sec:compression} should be noted. A signed basis $x = x^0 + Qz$
does not automatically have the convex-combination property, because entries
of $Q$ may be negative; convexity of the feasible region is then maintained
only by imposing the explicit linear system $x^0 + Qz \ge 0$, $A(x^0 + Qz) =
d$, and the linear minimization oracle is no longer the standard OD simplex
oracle. For the column-based FW extension the cleanest division of labor is
therefore: use the (flow-weighted) SVD to \emph{discover} latent structure and
candidate groupings, construct nonnegative convex route-family columns from
that structure, and optimize the family flows over OD simplices. This is the
complementarity already anticipated in Section~\ref{sec:grouped}, sharpened
here into an exact statement.

\subsection{Recorded Computational Checks}
\label{sec:fw-origin-evidence}

The reformulation is exercised in a separate reproducibility
package\footnote{\texttt{compressed\_TAP\_repro\_package\_v1\_1\_FW\_origin},
directory \texttt{fw\_origin\_extension/}; the numbers below trace to
\texttt{results/reference\_sioux/FW\_controlled.csv},
\texttt{ORIGIN\_controlled.csv}, and \texttt{cpp\_kernel\_*.csv}. Wall-clock
ratios are hardware-dependent (recorded single-thread runs); objective gaps,
atom counts, storage reductions, and demand residuals are deterministic.}
implementing FW and an origin-block conditional gradient scheme directly in
the compressed coordinates \eqref{eq:compressed-fw}, with exact BPR line
search, a common $10^{-5}$ relative-gap stopping rule, and exact demand
feasibility at every iterate (recorded demand residuals below
$1.3 \times 10^{-11}$). Table~\ref{tab:fw-origin-grouped} reports the
controlled route-rich experiment: $20$ OD blocks whose admissible route pool
grows with $K$ while the grouped representation holds $100$ atoms fixed.

\begin{table}[htbp]
\centering
\caption{Controlled route-rich experiments: full path pool versus grouped
convex route families (four families and one explicit major per OD block).
Objective gaps are relative to the full-pool solution of the same method;
speedups are median solve-time ratios of the recorded single-thread runs.}
\label{tab:fw-origin-grouped}
\begin{tabular}{lrrrrrr}
\toprule
Method & $K$ & Full atoms & Grouped atoms & Obj.\ gap (\%) & Speedup & Storage \\
\midrule
Frank--Wolfe & 8  & 160 & 100 & 0.190 & $1.18\times$ & $1.6\times$ \\
Frank--Wolfe & 16 & 320 & 100 & 0.114 & $2.10\times$ & $3.2\times$ \\
Frank--Wolfe & 32 & 640 & 100 & 0.045 & $5.17\times$ & $6.4\times$ \\
Origin-block & 8  & 160 & 100 & 0.190 & $1.14\times$ & $1.6\times$ \\
Origin-block & 16 & 320 & 100 & 0.114 & $1.73\times$ & $3.2\times$ \\
Origin-block & 32 & 640 & 100 & 0.044 & $5.21\times$ & $6.4\times$ \\
\bottomrule
\end{tabular}
\end{table}

Two further recorded checks connect this section to the weighted-SVD analysis
of the earlier sections. First, the flow-weighted construction transfers to
standard FW as a feasible \emph{pool-screening} mechanism: weighted pivoted QR
selects actual paths (no signed vector enters the simplex), after which
ordinary restricted-pool FW is run. At $K = 16$ the screened pool attains a
$0.015\%$ objective gap versus $0.308\%$ for unweighted screening (speedup
$1.53\times$ versus full FW); at $K = 32$, $0.020\%$ versus $0.245\%$
($1.12\times$). Second, a standalone C++17 kernel benchmark of the compressed
forward/adjoint/linear-minimization operations on the packaged real-network
path sets (Sioux Falls and Chicago Sketch at $K = 32$; $96$ full rows versus
$15$ grouped rows) records single-thread kernel speedups of $3.10\times$ and
$1.68\times$, respectively.

These checks are deliberately modest in scope: they certify feasibility,
exactness of conservation, and the predicted fixed-dimension behavior (the
grouped variable count stays constant while the full pool grows with $K$, and
the grouped model also requires fewer FW iterations in the richer cases),
rather than making large-network solver claims, which belong to the protocols
of Section~\ref{sec:experiments}.

\subsection{Solver Portability of the Latent Representation}
\label{sec:solver-portability}

The convexity-preservation argument implies a testable prediction: the
compression gain should be a property of the \emph{representation layer}, not
of any particular solver. Because problem \eqref{eq:compressed-fw} is a smooth
convex program over a product of simplices, any method for this constraint
geometry applies. A dedicated study\footnote{%
\texttt{compressed\_TAP\_repro\_package\_v3\_1\_LPG}, directory
\texttt{repro/fw/}; numbers trace to
\texttt{results/solver\_portability/solver\_portability.csv} and the recorded
run of \texttt{repro/interface/integrated\_check.py}. Wall-clock ratios are
hardware-dependent single-thread recordings; objective values, representation
gaps, and residuals are deterministic.} tests this by comparing the two
solvers of Section~\ref{sec:fw-interpretation} --- L\,$+$\,FW and
L\,$+$\,PG --- which differ in exactly one component, the update operator,
while sharing the atom set, the latent gradient, the OD demand simplices, the
objective, and the duality-gap stopping certificate. The protocol enforces
matched accuracy per representation: the restricted optimum $f^*_\rho$ of each
representation is first computed to duality gap $10^{-10}$, each timed run
stops at first crossing of $f^*_\rho(1 + 10^{-6})$, speedups are
\emph{within-solver} ratios (explicit-pool time over compressed time for the
same method), and the representation error $(f^*_\rho - f^*_E)/f^*_E$ is
reported separately, never folded into a speedup.

\begin{table}[htbp]
\centering
\caption{Solver portability, controlled sweep ($4$ OD pairs, $8$ latent atoms
per OD plus one explicit major in RC, growing route pool $K$).
Matched-accuracy within-solver speedups of the compressed representation (RC)
over the explicit pool (E), and representation gaps of RC and of pure latent
C (no explicit major). Recorded single-thread runs.}
\label{tab:solver-portability}
\begin{tabular}{rrrrr}
\toprule
$K$ & L\,$+$\,PG speedup & L\,$+$\,FW speedup & RC gap (\%) & C gap (\%) \\
\midrule
8   & $1.02\times$ & $1.00\times$ & 0.00 & 0.00 \\
32  & $1.17\times$ & $1.07\times$ & 1.19 & 2.89 \\
128 & $3.00\times$ & $0.90\times$ & 1.78 & 6.26 \\
512 & $7.70\times$ & $2.46\times$ & 1.85 & 6.23 \\
\bottomrule
\end{tabular}
\end{table}

Three observations from Table~\ref{tab:solver-portability} and the
accompanying real-network runs deserve emphasis. First, the projected-gradient
speedup grows monotonically with the pool-to-atom ratio, reaching
$7.70\times$ at $K = 512$: PG's per-iteration cost (link loading plus an
$O(K \log K)$ projection) scales with the representation while its iteration
count does not, which is exactly the dimension-dependence the portability
argument requires. Compression benefits projection-based methods at least as
much as linear-oracle methods. Second, the FW wall-time caveat is real: in
this Python implementation FW's time is iteration-count bound (the $O(1/k)$
tail chasing a $10^{-6}$ target) with size-independent per-iteration overhead,
so its wall-time speedup stays near one until $K = 512$; FW's per-iteration
dimension gains are established separately at kernel level
(Section~\ref{sec:fw-origin-evidence}). Wall-time and kernel evidence must not
be conflated. Third, the support mechanism of the earlier sections reappears
at its most dramatic on the congested Chicago Sketch instance ($K = 32$,
doubled demand): pure latent C --- convex atoms with nominal shares and no
explicit major --- fails by $149.8\%$, because nominal-share convex
combinations cannot concentrate mass on the dominant path; a single explicit
major per OD (RC) repairs this to $0.53\%$, with within-solver E/RC speedups
of $95\times$ (L\,$+$\,FW) and $4.3\times$ (L\,$+$\,PG) on that instance. On the
anchor-supported Sioux Falls instance all representations are exact and both
solvers converge essentially immediately.

Finally, an integrated five-solver acceptance check runs one exported
problem-layer instance through SLSQP on the anchor-eliminated grouped decoder,
the C++ projected-gradient kernel, FW on the interface atoms, FW on the full
path pool, and L\,$+$\,PG on the same interface atoms, and requires agreement of
their optima. On the spread-demand Chicago instance, L\,$+$\,PG matches the SLSQP
reference to $3.7 \times 10^{-14}$ relative --- on the identical atoms where
vanilla FW's sublinear tail stops at $9.8 \times 10^{-5}$. The verified
equivalence chain thus reads: ALM $=$ simplex projected gradient $=$
reduced gradient $=$ L\,$+$\,FW $=$ L\,$+$\,PG, all through the standard problem-layer
interface, certified by a common duality gap. This is the operational content
of Proposition~\ref{prop:convexity-preservation}: one convex representation,
many interchangeable solvers.

\subsection{Beyond Traffic: Major--Latent Prototypes for General Column Pools}
\label{sec:beyond-traffic}

Nothing in Proposition~\ref{prop:convexity-preservation} is specific to
traffic assignment: the argument uses only that the decision vector is a
nonnegative combination of columns, that the observable quantity is linear in
the decisions, and that the objective is convex in the observable. Two small
prototypes\footnote{\texttt{major\_latent\_prototypes/} (LP-relaxation
prototypes, \texttt{examples/run\_all.py}); deterministic column counts and
relative errors verified by a fresh run on 2026-07-18. The prototypes
deliberately exclude branch-and-price, integer recovery, and software
integration; they test only the representation layer.} exercise the same
representation, $x = Wz$ with effective matrix $\widetilde B = BW$ and cost
$\widetilde c = W' c$ (the operator $W$ is the general-column-pool analogue of
$G$ in \eqref{eq:G-reconstruction}), on classical column-pool linear programs:
cutting stock, where the original columns are cutting patterns, and unit
commitment, where the original columns are feasible generator schedules. Each
prototype compares the full explicit pool (E), fully latent convex atoms with
no explicit major columns (C), and explicit majors plus latent minor atoms
(RC).

\begin{table}[htbp]
\centering
\caption{Beyond-traffic prototypes (LP relaxations): column counts and
relative objective error of the compressed representations versus the full
explicit pool.}
\label{tab:beyond-traffic}
\begin{tabular}{lrrrrr}
\toprule
Problem & E cols & C cols & RC cols & C error (\%) & RC error (\%) \\
\midrule
Cutting stock   & 10  & 5  & 8  & 4.51 & 0.00 \\
Unit commitment & 162 & 18 & 24 & 5.60 & 0.38 \\
\bottomrule
\end{tabular}
\end{table}

The pattern of Table~\ref{tab:beyond-traffic} reproduces, in miniature, the
support mechanism of Section~\ref{sec:solver-portability}: pure latent
compression (C) incurs a visible objective error because frozen convex shares
cannot concentrate weight on the dominant columns, while retaining a few
explicit major columns (RC) recovers the optimum to $0.00\%$ and $0.38\%$ at
$5\times$ and $7\times$ column reduction, respectively. Feasibility residuals
are at machine precision in all cases, as the convexity argument requires.
These prototypes support the broader claim of the concluding section: the
major--latent construction is a representation-layer device for column-pool
convex programs in general, with traffic assignment as its most developed
instance.

\subsection{Summary Statement}
\label{sec:fw-origin-summary}

The latent route-family representation is convexity preserving. Each latent
column is constructed as a nonnegative convex combination of feasible path
columns within the same OD pair. The compressed decision variables remain on
the OD demand simplex, and their linear reconstruction produces nonnegative
path flows satisfying exact OD conservation. Consequently, the represented
feasible set is a convex subset of the original path-flow polytope. Because
link flows are linear in the compressed variables and the Beckmann potential
is convex under nondecreasing link costs, the resulting compressed assignment
remains a convex program. Standard Frank--Wolfe updates retain their classical
interpretation as convex combinations of feasible assignments. This is a
central methodological property of the framework, not an implementation
detail: it is what allows the compression to be embedded in FW, origin-based,
and column-generation solvers without altering their convergence theory.
